\section{Building a quality database from the Web}
\newenvironment{tightenumerate}{\begin{list}{\labelenumi}{\usecounter{enumi}%
      \setlength{\topsep}{0.05in}%
      \setlength{\itemsep}{0.02in}%
      \setlength{\parsep}{0in}
      }}{\end{list}}


The rich and diverse information available on the Web lends us enormous
opportunity to build a comprehensive and high-quality 
database for cross-modality search. At the same time,
the very richness and diversity
poses many interesting challenges that need to be addressed.
For example, given the vast size of the Web, how can we 
``identify'' and ``extract'' relevant 
data from the Web automatically and add them to our database? 
Also, how can we {\em integrate} {\em heterogeneous} data on the Web
when they are based on different ``vocabulary'' or ``schema''?
In the following subsections,
we will describe some of the challenges that we plan
to address under this data-collection thrust.

\subsection{Automatic data collection}\label{sec:sampling}
While there exist an enormous amount of information available on the Web,
the information is often {\em scattered} and {\em unorganized}
due to the lack of a central authority. Any person can publish
any information at a Web site that he chooses, so some information
might be ``buried'' at a remote site that few people are aware of.
Also, much of the information on the Web might not be very useful
for a particular purpose. For example, if we want
to build a database of artistic paintings,
the pages related to music will not be very
useful for our database. Given these facts, our data collection
should achieve the following two goals:

\begin{itemize}
\item {\em Cover a wide range of Web:}
  As we argued, information on the Web is quite unorganized and dispersed,
  so we need to explore different parts of the
  Web in order to collect most of the potentially useful data.
  Unless we visit pages (or multimedia objects) and check their 
  relevance, it is difficult to tell whether they are useful or not.
\item {\em Focus on the relevant parts of the Web:}
  Given the size and the growth rate of the Web, it is difficult 
  to download every page (or object) on the Web. Clearly, 
  our data collection should be focused on the part of the 
  Web {\em relevant} to our task.
\end{itemize}

The goals stated above seemingly conflict: We want to explore
many different parts of the Web, while focusing only on
the relevant parts. As a potential solution
to the conflicting goals, we plan to study a {\em sampling-based approach}.
Under a sampling-based data collection, we download
pages using the following algorithm.

\begin{tightenumerate}
\item Split the Web into $N$ partitions.
\item From a partition $p_i$, sample $k$ pages.
\item If many of the $k$ sample pages are ``relevant,''
      download all pages in the partition.
\item Otherwise, move to the next partition.
\end{tightenumerate}

To help our discussion, let us assume that we partition the Web
by {\em Web sites}, and we try to build a database of artistic paintings.
For this scenario, the above algorithm would download, say, $100$ 
pages and/or images from a Web site and check their relevance 
to artistic paintings. If many of them are relevant, it downloads more pages 
from the site. Otherwise, it skips
to the next site. By collecting data this way, it can
potentially {\em explore} many different parts of the Web,
while {\em focusing} only on the relevant parts. It only commits
significant download resources to a site only when the site has
a high chance to have many relevant pages (or images).

Clearly, we need to address many open issues 
to make the above algorithm work. For example, how
should we partition the Web to make the algorithm most effective?
In the above example, we partitioned the Web based on Web sites,
but we need to carefully select the right partitioning function,
because it will affect performance very significantly:
If more ``related'' pages are put into the same partition, the
algorithm can collect many relevant pages. 
Also, how many samples should we
download from each partition? If we sample more pages, we can
judge the relevance of a partition with more confidence, but
we may waste our resources sampling potentially irrelevant pages.

\subsection{Automatic data extraction}
Once we collect data from the Web, we need to {\em extract}
relevant information from the collected data,
so that we can load the information into our database.
For example, if we want to build a database of artistic paintings,
we need to extract artistic painting from downloaded pages 
and store them in our database.

However, the data on the Web are often unstructured 
and follow different ``schemas,'' so automatic extraction
of information is a challenging task.
How can we tell whether a downloaded image is an artistic 
painting or not? Also, to support cross-modality search,
we need to associate relevant {\em keywords} to the extracted
images. How can we automatically figure out relevant keywords?
Note that most of the multimedia objects on the Web are embedded
on a Web page. Therefore, we may use the ``surrounding-text''
to each object to predict the {\em genre} of an object
and identify associated keywords. 

\begin{figure}
  \centering
  \subfigure{%
    \includegraphics[width=2in]{regpage1.eps}}%
  \qquad% 
  \subfigure{%
    \includegraphics[width=2in]{regpage2.eps}}%
  \caption{Pages from the Metropolitan Museum of Art Web site}
  \label{fig:regpage}
\end{figure}

\paragraph{Exploiting regularity}
Pages from the same Web site often follow the same format.
For example, Figure~\ref{fig:regpage} shows two pages from the Metropolitan
Museum of Art Web site. Note that both of the pages have the same title,
and the same side bars. Furthermore, 
the images and their explanations are arranged side-by-side,
in the same format on both pages. Therefore, if we can detect
this regularity and identify the pattern, we may automatically
extract images and their explanations. In general, we plan to
explore the following algorithm for automatic data extraction.
For concreteness, we assume that we are building an image database,
but note that the algorithm can be applied to any multimedia objects,
such as music and video.

\begin{tightenumerate}
\item From the Web pages that are relevant to the images that we want
  to collect (or from the database that we have already built),
  extract a set of relevant keywords.
\item Identify the Web pages that contain
  images and the relevant keywords in them.
\item From the pages, detect how the images and the keywords are arranged.
  \label{item:algo-pattern}
\item If more than $k$ pages have the same arrangement of images and
  keywords, extract images and keywords from the pages and load them
  to the database.
\end{tightenumerate}

Depending on the value of $k$, the algorithm will extract 
different quality of data.
If $k$ is large, the algorithm extracts images only when a large number
of pages have the same arrangement (i.e., when their is a high evidence), 
so the images have a high chance to be related to the extracted keywords.
On the contrary, if $k$ is small, the algorithm will extract more
images, but there is a higher chance that the images are not
actually related to the keywords.
Also, 
there exist a multitude of ways to represent the arrangement patterns
in step~\ref{item:algo-pattern} of the above algorithm.
For example, we may use a regular expression to represent the patterns
or we may construct a parse tree of the HTML page.
We plan to study these issues carefully to learn how we can extract
high quality data from the Web automatically.


\subsection{Efficient data maintenance}
Data on the Web are constantly changing: New pages
are constantly created, obsolete pages are continually deleted
and existing pages are updated over time.
For example, pages on the CNN Web site are updated
whenever there are new developments in the world,
and the IBM Web site is completely overhauled,
if the company goes through a major re-organization.
Because these changes
are made autonomously and the Web sites do not inform our database
of these changes, we have to periodically {\em poll} 
and {\em download} pages and/or multimedia objects
to check for changes and update our database.

Effective refresh of a local database introduces many
interesting challenges. For example, how often should we
refresh a page (or a multimedia object) to maintain it up-to-date?
How should we allocate our limited download resources to different
objects, when the objects are of different size and popularity?
In~\cite{cho00a}, we proposed two definitions for the ``freshness''
of a database and we addressed some of these questions
experimentally and theoretically. We plan to build on our previous work 
and extend it in the following ways.

\begin{itemize}

\begin{figure}
  \centering
  \includegraphics[width=2.5in]{mod-int-all.eps}
  \caption{Average change interval of Web pages}
  \label{fig:change-interval}
\end{figure}
\item {\bf Change frequency:}
Web pages are changing at vastly different rates:
Pages on the CNN Web site change at least once a day on average,
while pages on the IRS Web site may be updated only once a year.
Figure~\ref{fig:change-interval} 
shows an experimental results that we reported in \cite{cho00b}.
In \cite{cho00b}, we traced the change history of 720,000 Web pages for
4 months and estimated how often Web pages change on average. 
The horizontal axis in the graph shows the average change interval
of pages and the vertical axis shows the fraction of pages changed
at the given rate. The graph clearly shows that Web pages have
very different change frequencies.

We can exploit these different change frequencies of pages to improve
the freshness of the database. For example, assume that we 
have guessed that the pages on the CNN Web site changes more
often than the pages on the IRS Web site after we visit the
sites several times. Then we may want to visit the CNN Web site
more often than the IRS Web site. 

In \cite{cho01a}, we studied how we can estimate
the change frequencies of Web pages based on its past change history,
and in \cite{cho01a} we studied how we should revisit pages once we
know their change frequencies. However, in the previous studies
we assumed that we first visit pages for a certain period of time
to {\em estimate} page change frequencies, and after 
this period, we use the estimate to adjust page revisit frequencies.
That is, we assumed a {\em static} scheme where we estimate change
frequencies only once and use the estimates for the rest.
We plan to extend this framework into a {\em dynamic} scheme,
where we can {\em continuously} and {\em dynamically} estimate 
change frequencies and adjust revisit frequencies appropriately. 
For example, the change
frequency of a page itself may change over time, so we may need
to re-estimate its change frequency and adjust our policy.
Also, our initial guess for the change frequency may turn out to be wrong
and we may want to adjust its revisit frequency.

\item {\bf Correlated updates:}
In certain cases, the changes to some pages are highly correlated,
and we may exploit this correlation to improve the ``freshness''
of our database.
For example, consider the pages on the Linux Operating System.
When a new major release of the Linux kernel is announced, 
many of the Linux-related
pages would change to reflect this fact. Therefore, if we notice
some Linux pages have changed, we may wish to refresh {\em all other} 
Linux-related pages, expecting that they may also have changed. 
We can formalize this idea as follows:

\begin{tightenumerate}
\item Find pages related to particular topics, say the Linux Operating System.
\item From each set of pages, sample $k$ pages.
\item If many of the $k$ sample pages have changed from the previous download,
      download all pages in the set.
\item Otherwise, move to the next set of pages.
\end{tightenumerate}

Notice the similarity of the above algorithm to the one in 
Section~\ref{sec:sampling}. In fact, we may generalize the above algorithm,
so that it first partitions the Web in to $N$ partitions and sample
$k$ pages from each partition to check for changes.
We may partition the Web based the contents of a page, the link
or even simple directory (or domain) structure of the Web.
Clearly, we will have a similar set of issues that need to be addressed.
For example, what will be the best partitioning method that can put
corrrelated pages in the same partition?
How many samples should we take from each partition? 
Should we sample different numbers of pages depending on the size and the
importance of a partition?

Also, what will be a good statistical model to capture the correlation
among pages? A good model will be important to compare 
various refresh algorithms and to understand why the
algorithms give certain results. We will explore these questions 
experimentally using real data
and we will rigorously study the experimental results through
theoretical analysis.

\item {\bf Pattern detection:}
Not all changes are equal.
For example, a page on a commercial Web site may change very
often, but the change might be only at its advertising banners.
In contrast, a page on the IRS Web site may change only once a year,
but the pages may be completely updated. 
We may improve the freshness of our database by considering
different patterns of changes and adjusting refresh policies properly.
For example, we may want to wait for a while for the pages
that change often yet slightly, in order to collect a set of changes 
in a batch.
Roughly, we may categorize change patterns into four classes:
1) changes occur only at a particular part of a page and the 
changes are ``important'' to the users 2) changes occur only at one part,
but the changes are not important 3) changes occur only at a small part
of a page in each update, but the location of the change move 
in each update 4)
the page is completely overhauled at every update. Depending
on which change pattern a page follows, we may want to adop different
refresh policies.

To realize this idea, we need to address many interesting questions:
How can we detect different patterns of changes?
We may not be able to store all past versions of a page due to 
our limited storage space, and detecting change patters in this
scenario will be difficult.
Also, exactly how should we adjust revisit policies 
based on the detected patterns?
We need to define the clear notion of ``freshness'' and 
compare tradeoffs carefully to answer these questions.
\end{itemize}


\subsection{Heterogeneous data integration}
  

\subsection{Related work}
There exists a large body of literature that studies how to automatically
collect data from the Web. This data collection is often conducted by
a program, called a {\em Web crawler}, which traverses the link structure
of the Web to discover and download pages. 

\paragraph{Focused crawling}
Because a crawler has limited download and storage resources,
a {\em focused crawler}
tries to collect only the pages related to a 
{\em particular topic}~\cite{chak99a,chak99b,dil00,mukh00,agga01,
menc01,cho98,mcca99,hers98}. 
In most of the existing work,
researchers have studied how to analyze keywords and the link
structure around a page to judge its relevance to a topic.
For example, in our previous work~\cite{cho98},
we found that a page has a high chance 
to be related to a certain topic, when the page is linked 
to a page that is already found to be relevant.
References~\cite{chak99a,chak99b,dil00} and~\cite{agga01} 
propose more improved versions of this basic idea. 
The crawler described in \cite{chak99a,chak99b} has two modules,
a classifier and a distiller, that judge the relevance of a page 
and determines page download order. The decision is 
based on the relevance of the page and its potential proximity to 
other relevant pages. 
Researchers have also tried to use {\em machine learning} techniques
to build a focused crawler. For instance,
reference~\cite{mcca99} adopts {\em reinforcement learning} 
and applies Bayes classifiers to the full text and anchor text of a page 
to classify a page.

We believe our {\em sampling-based} approach 
has a great potential to improve the effectiveness of a crawler even
further. Since it {\em samples} various parts of the Web,
it has a less chance to be {\em trapped} in a small cluster of relevant
pages and miss other more relevant pages. Also, the overhead from sampling
can be minimal, because we use only a small portion of our download
resources for sampling. Also note that
a {\em sampling} crawler can still incorporate many of the mechanisms 
proposed in the literature. For example, we can use the Bayes classifier
to judge the relevance of a page, potentially improving the accuracy of 
judgement.


\paragraph{Data refresh}
A lot of work has been done on maintainging 
replicated data consistent~\cite{bern84,alon90,
bern80,barb92,carv82,colby97,krish91,krish94,ladin92,olston00}. 
Most of the work studies the 
tradeoff between consistency and read/write 
performance~\cite{krish94,ladin92,alon90,olston00}
and tries to guarantee a certain level of consistency
under various definitions of consistency~\cite{alon90,barb92,bern80}. 
However, most of the work assumes a {\em push} model,
where the updates to the data sources are {\em not} autonomous 
and the sources notify the replicated data of the updates. 
In contrast, we assume a {\em polling} model where
updates are made {\em independently} and {\em autonomously} at the sources.

References~\cite{coff97,cho00a} study how a crawler can maintain its index
up-to-date when the Web pages are independently updated. While 
both work shows optimal refresh strategies when the crawler know how often
pages change, neither of them addresses how the crawler can estimate the change
frequencies of pages. 
References~\cite{wolman99,evol-att,evol-cache,BC2000,cho00b} 
experimentally estimate change frequencies
of Web pages. In~\cite{cho01a}, we studied change frequency estimation
when a crawler has limited change history of a page.
However, we mainly separated frequency estimation 
from actual page refresh, and we assumed that the estimation is done
only once. We plan to extend the previous work to {\em combine} 
page estimation with page refresh, where we estimate change frequencies
dynamically and continally over time.

\paragraph{Automatic data extraction}
%Data extraction from Web pages has got much attention recently.
%References~\cite{cho97}

\paragraph{Heterogeneous data integration}

%%% Local Variables: 
%%% mode: latex
%%% TeX-master: "main"
%%% End: 
